% Long caption; ordinary float; the caption is taller than the artwork.
\begin{figure}[htbp]
\centering
\rule{180pt}{45pt}
\caption[Long caption test]{A long caption must remain in the outer margin
even when it is taller than the exhibit. The two columns flow independently.
It must not silently return this paragraph to the main reading
column, shrink its type, crop its final sentence, or overlap the paragraph that
follows. This deliberately verbose test caption exceeds the former quarter-page
limit. Its wording repeats the constraints so the test measures real wrapped
prose rather than an artificial fixed-height box. The final sentence must be
visible on the same page as the black rectangle, and the caption must remain
on the right-hand side of this odd-numbered page.}
\label{fixture:long}
\end{figure}
This paragraph must begin below the rectangle, alongside the tall caption.
\marginnote{A second margin note must clear the caption without moving body text.}
\clearpage
% Caption before a table, nested inside center.
\begin{table}[htbp]
\begin{center}
\caption{The caption belongs to a table even inside another environment.}
\label{fixture:nested-table}
\begin{tabular}{ll}
\toprule Choice & Result\\\midrule Left & Outer margin\\\bottomrule
\end{tabular}
\end{center}
\end{table}
\clearpage
% Source encountered on an odd page; shipped on the following even page.
\noindent Source page.\par\vspace*{400pt}
\begin{figure}[t]
\centering\rule{180pt}{220pt}
\caption{A moved float takes its caption with it, onto the next even page.}
\label{fixture:moved}
\end{figure}
\clearpage
% A figure* is still composed within the Book column, beside its caption.
\begin{figure*}[htbp]
\centering
\begin{tikzpicture}
\draw (0,0) rectangle (10.5,2);
\end{tikzpicture}
\caption{A starred figure keeps its caption beside its top, never below it.}
\label{fixture:wide}
\end{figure*}
\clearpage
% Starred captions are also marginal but do not increment the counter.
\begin{figure}[htbp]
\centering\rule{120pt}{40pt}
\caption*{An unnumbered caption also belongs in the outer margin.}
\end{figure}
\clearpage
\noindent A page-breaking table is encountered late on this odd page.
\par\vspace*{380pt}
\begin{xltabular}{\textwidth}{@{}lX@{}}
\caption{A page-breaking table keeps its caption in the same outer margin.
This caption is deliberately longer than the space left on its source page.
The build measures its height and reserves that much space before starting
the table on the next page. It must not cross the running foot, disappear,
or become an ordinary heading above the rows. The column headings repeat
when the table continues, but this caption and its number appear only once.
Its last sentence must remain visible next to the first part of the table,
on the left side of the even page that actually receives it.}
\label{fixture:longtable}\\
\toprule Choice & Result\\\midrule
\endfirsthead
\toprule Choice & Result\\\midrule
\endhead
\bottomrule\endlastfoot
A & \rule{0pt}{80pt}The table rows remain in the main column.\\
B & \rule{0pt}{80pt}The caption is not a heading above the table.\\
C & \rule{0pt}{80pt}Its number and list entry are still generated once.\\
D & \rule{0pt}{80pt}A tall row tests the page break.\\
E & \rule{0pt}{80pt}Another row preserves its table heading.\\
F & \rule{0pt}{80pt}No caption is repeated on the continuation.\\
G & \rule{0pt}{80pt}The last page is still part of the same table.\\
H & \rule{0pt}{80pt}The final row remains visible.\\
\end{xltabular}
\clearpage
\begin{lstlisting}[caption={A short code listing still reserves its entire margin caption, even when the description is taller than its two lines of code.},label={fixture:listing}]
acquire(file);
commit();
\end{lstlisting}
This paragraph follows the listing, alongside its caption.
\clearpage
\begin{figure}[htbp]
\centering
\begin{minipage}{.92\textwidth}
\begin{lstlisting}[caption={A listing inside a narrow minipage and a float uses the page margin, not the minipage edge.},label={fixture:nested-listing}]
release(file);
\end{lstlisting}
\end{minipage}
\end{figure}
\clearpage
% A portrait-shaped margin object anchored near the foot must be raised as
% a whole, including its description. Never crop it or push body text away.
\noindent Portrait regression page.\par\vspace*{470pt}
\noindent The portrait belongs with this sentence.\marginnote{%
  \rule{\marginparwidth}{145pt}\par A complete portrait and its explanation.}
\par Body text continues normally.
\clearpage
\begin{pdsession}[The first client holds the lock; the second must wait.]
$ client-a acquire file.txt
acquired
$ client-b acquire file.txt
held by client-a
\end{pdsession}
The terminal explanation is in the margin, not above the transcript.
\clearpage
% A deferred medium-height exhibit should share its next page with prose,
% rather than qualify for an eager half-full float-only page.
\noindent\rule{0pt}{290pt}Source text before a deferred exhibit.\par
\begin{figure}[!htbp]
\centering\rule{180pt}{350pt}
\caption{A medium-height exhibit shares its page with continuing prose.}
\label{fixture:shared-float-page}
\end{figure}
\newcount\pdflowtestcount\pdflowtestcount=0
\loop
  \advance\pdflowtestcount by1
  \noindent Flow continuation. This ordinary paragraph is deliberately long
  enough to keep the main text flowing around the deferred exhibit. It must
  share the exhibit's page when space remains below the artwork, without
  changing the caption size or borrowing its margin.\par
\ifnum\pdflowtestcount<14\repeat
\clearpage
% A short table fits below this paragraph; the old 40%-page fallback did not.
\noindent Short xltabular source.\par\vspace*{390pt}
\begin{xltabular}{\textwidth}{@{}lX@{}}
\caption{A short measured caption stays beside its table.}
\label{fixture:short-xltabular}\\
\toprule Field & Value\\\midrule\endfirsthead
\toprule Field & Value\\\midrule\endhead
\bottomrule\endlastfoot
Probe & Short xltabular row.\\
\end{xltabular}
\clearpage
\noindent Short longtable source.\par\vspace*{390pt}
\begin{longtable}{@{}ll@{}}
\caption{A directly authored longtable uses the same measured reservation.}
\label{fixture:short-direct-longtable}\\
\toprule Field & Value\\\midrule\endfirsthead
\toprule Field & Value\\\midrule\endhead
\bottomrule\endlastfoot
Probe & Short longtable row.\\
\end{longtable}
\clearpage
% A framed passage must split without mdframed's retry loop exhausting.
\noindent Frame split source.\par\vspace*{390pt}
\begin{pdexample}{A breakable passage}
\pdflowtestcount=0
\loop
  \advance\pdflowtestcount by1
  A worked example may continue over a page turn. Its full text must remain
  visible, with readable line spacing and no oversized frame crossing the
  footer. Strong paragraph preferences must not make the frame unsplittable.
\ifnum\pdflowtestcount<12\repeat
Frame split final sentence.
\end{pdexample}
\clearpage
\noindent Short frame source.\par\vspace*{500pt}
\begin{pdexample}{A short passage with a heading that wraps onto a second line}
The full paragraph follows its heading onto a fresh page when there is not
enough room to begin safely. It must not enter a failed split retry.
Short frame final sentence.
\end{pdexample}
\clearpage
\noindent Protocol split source.\par\vspace*{390pt}
\begin{stpprotocol}[Page-breaking protocol]
\label{fixture:split-protocol}
\begin{enumerate}
\pdflowtestcount=0
\loop
  \advance\pdflowtestcount by1
  \item A protocol step explains the action and the conditions that make it
  valid. A long protocol continues with the following steps on the next page,
  rather than moving its whole rectangle and leaving this space empty.
\ifnum\pdflowtestcount<8\repeat
\end{enumerate}
Protocol split final sentence.
\end{stpprotocol}
\clearpage
\noindent Heading and example source.\par\vspace*{460pt}
\subsection{Framed subsection}
\begin{pdexample}{An example directly after its heading}
First framed sentence. The heading and the start of this example must
remain together, even when the example eventually continues on another page.
\end{pdexample}
\endinput
